\subsection{CP-07: reintentos con retroceso exponencial}
\label{sec:cp07}

\FichaCaso
  {CP-07}
  {David Estevan Rodriguez Jurado}
  {09/11/2026}
  {Reintentos con retroceso exponencial}
  {Joan Emmanuel Orduz Chia}
  {11/11/2026}
  {Comprobar el atributo de calidad de confiabilidad y resiliencia del SDK: las operaciones idempotentes reintentan solo fallos transitorios, con retroceso exponencial y un tope, y las que no son idempotentes no se reintentan para no cobrar dos veces. El caso verifica los tiempos, el número de intentos, el error final con su causa y el comportamiento de \texttt{createPayment()} frente a una pasarela que falla y luego responde.}
  {No funcional (confiabilidad). No es un CRUD ni una prueba de seguridad.}
  {\textbullet\ SDK compilado y dependencias de \path{examples/} instaladas.\newline
   \textbullet\ API de Simulación arriba en \texttt{http://localhost:3000} para los escenarios de integración.\newline
   \textbullet\ Las suites de timeout y de errores HTTP del SDK levantan su propio servidor en \texttt{127.0.0.1} con un puerto efímero y no necesitan el simulador.\newline
   \textbullet\ Sin credenciales reales.}
  {Local. Node.js 20. Las suites del SDK corren contra un servidor HTTP local en un puerto efímero; los escenarios corren contra la API de Simulación en \texttt{http://localhost:3000}. No se usa red externa ni credenciales de sandbox.}

\subsubsection{Objetivo}

Comprobar que el SDK distingue un fallo transitorio de uno de negocio, que la espera entre intentos crece de forma exponencial hasta un tope, que el número de intentos es el configurado y que el error final conserva la causa. Comprobar además que \texttt{createPayment()} no reintenta: un cobro que falla dos veces y prosperaría en el tercer intento debe devolver el error la primera vez, no cobrar en silencio.

\subsubsection{Procedimiento}

\begin{TablaPasos}{Pasos del caso CP-07}
  1 &
  Ejecutar las suites de reintentos, timeout y errores HTTP. &
  \texttt{cd sdk \&\& npx jest RetryHandler} \path{KitPagos.timeout} \path{KitPagos.http-errors} \texttt{--coverage=false --verbose}. &
  Los 3 archivos y sus 55 pruebas pasan, 0 fallos. Con \texttt{jitterMs:0}, \texttt{calculateDelay(0..2)} da 1000, 2000 y 4000, y el tope de 4000 se repite en los intentos siguientes. &
  55 de 55. & Pasó & Suite \\ \hline
  2 &
  Verificar el número de intentos y las pausas. &
  La prueba \emph{retries up to the configured limit before succeeding on the last attempt} de \path{RetryHandler.test.ts}. &
  La operación se llama 4 veces (1 intento inicial y 3 reintentos), \texttt{sleep} se llama 3 veces y las pausas son \texttt{[1000, 2000, 4000]}. Solo se reintentan los códigos transitorios. &
  Pasa la prueba. & Pasó & Suite \\ \hline
  3 &
  Verificar que un error de negocio no se reintenta. &
  La prueba \emph{rethrows non-transient errors immediately without retrying or pausing} de \path{RetryHandler.test.ts}. &
  \texttt{INVALID\_CREDENTIALS} se relanza en el primer intento, la operación se llama una vez y no hay pausa. &
  Pasa la prueba. & Pasó & Suite \\ \hline
  4 &
  Verificar el error final y su causa. &
  La prueba \emph{exhausts the retries and throws MAX\_\allowbreak RETRIES\_\allowbreak EXCEEDED with the last error as cause} de \path{RetryHandler.test.ts}. &
  Se lanza \texttt{MAX\_RETRIES\_EXCEEDED}; el mensaje dice \texttt{4 attempts (1 initial + 3 retries)}; \texttt{cause} es el último error y \texttt{originalPayload} se conserva. &
  Pasa la prueba. & Pasó & Suite \\ \hline
  5 &
  Verificar el límite por petición y su reintento. &
  \path{KitPagos.timeout.test.ts}. &
  \texttt{getPaymentStatus()} corta a los 30~s por omisión con \texttt{GATEWAY\_TIMEOUT}; con un reintento termina en \texttt{MAX\_RETRIES\_EXCEEDED} y la causa es \texttt{GATEWAY\_TIMEOUT}. El servidor recibe exactamente 2 peticiones. &
  Pasan las pruebas. & Pasó & Suite \\ \hline
  6 &
  Levantar el simulador y ejecutar el recorrido de escenarios. &
  \texttt{cd simulator-api \&\& npm run dev}; en otra terminal, \texttt{cd examples \&\& npm run} \path{simulate:scenarios}. &
  Cada línea del recorrido imprime \texttt{OK} y el proceso termina con código 0. Entre otras filas: consulta \texttt{10602} se recupera tras dos fallos, consulta \texttt{10600} termina en \texttt{MAX\_RETRIES\_EXCEEDED}, consulta \texttt{10604} termina en \texttt{GATEWAY\_TIMEOUT} sin reintentos, y la lista de bancos con \texttt{sim\_flapping} se recupera. &
  61 de 61 en \texttt{OK}. & Pasó & Corrida 1 \\ \hline
  7 &
  Verificar que la creación no se reintenta. &
  \texttt{cd simulator-api \&\& npx jest} \path{scenario-through-sdk} \texttt{--coverage=false --verbose}. &
  La suite pasa sus 147 pruebas, entre ellas la del monto \texttt{10002} (\emph{FLAPPING, creation is not retried}): \texttt{POST /v1/api/payments} responde 502 con \texttt{GATEWAY\_SERVER\_ERROR}. Si el SDK reintentara, el tercer intento crearía el cobro y la respuesta sería 2xx. &
  147 de 147. & Pasó & Suite \\ \hline
  8 &
  Confirmar la asimetría en el código de la fachada. &
  Revisión de \path{sdk/src/infrastructure/facade/KitPagos.ts}. &
  \texttt{getPaymentStatus()} y \texttt{getPseBanks()} construyen un \texttt{RetryHandler}; \texttt{createPayment()} llama al adaptador directamente. La decisión está documentada para evitar un doble débito. &
  Como se espera. & Pasó & Código revisado \\ \hline
  9 &
  Repetir el recorrido para descartar una recuperación accidental. &
  El mismo comando del paso 6. &
  La segunda corrida imprime las mismas filas \texttt{OK} y vuelve a terminar con código 0; los contadores de fallo del simulador se reinician al recuperarse. &
  61 de 61 en \texttt{OK}. & Pasó & Corrida 2 \\ \hline
\end{xltabular}
\end{TablaPasos}

\subsubsection{Criterio de aceptación}

\begin{enumerate}
  \item Las 55 pruebas de \path{RetryHandler.test.ts}, \path{KitPagos.timeout.test.ts} y \path{KitPagos.http-errors.test.ts} pasan sin fallos.
  \item La espera es exponencial con tope: 1000, 2000 y 4000~ms sin jitter, y nunca supera los 4000~ms más el jitter configurado.
  \item Una operación que falla tres veces se intenta cuatro veces; un error no transitorio se relanza en el primer intento sin pausa.
  \item Al agotar los reintentos se lanza \texttt{MAX\_RETRIES\_EXCEEDED} con el último error como \texttt{cause}.
  \item El recorrido de escenarios termina con código 0 y cubre recuperación, agotamiento, timeout y lista de bancos.
  \item La creación con el monto \texttt{10002} devuelve 502 \texttt{GATEWAY\_SERVER\_ERROR} en el primer intento y no crea un cobro.
  \item \texttt{createPayment()} no usa \texttt{RetryHandler}; \texttt{getPaymentStatus()} y \texttt{getPseBanks()} sí.
\end{enumerate}

\subsubsection{Observaciones de diseño}

El caso separa dos niveles: las pruebas unitarias fijan la política con un reloj inyectado y sin red, y las pruebas de integración usan los montos reservados del simulador, que producen fallos deterministas (\texttt{10602} falla dos veces y se recupera, \texttt{10600} responde 500 y \texttt{10604} produce timeout). El monto \texttt{10002} es el que distingue una creación idempotente de una que no lo es: falla dos veces y prospera en el tercer intento, así que una respuesta 502 prueba que \texttt{createPayment()} no reintentó ni duplicó el cobro.

\subsubsection{Trazabilidad}

RF-05 (errores esperados), RF-07 (reintentos idempotentes), RF-10 y RF-11 (escenarios y errores del simulador); HU-05, HU-07 y HU-11.
